% Generated draft: requires mathematical and source review.
\documentclass[11pt]{article}
\usepackage[margin=25mm]{geometry}
\usepackage{fontspec}
\setmainfont{Latin Modern Roman}
\usepackage{amsmath,amssymb,mathtools,unicode-math}
\setmathfont{Latin Modern Math}
\usepackage{enumitem,needspace,xurl,hyperref}
\setlength{\parindent}{0pt}
\setlength{\parskip}{5pt}
\setlength{\emergencystretch}{3em}
\newcommand{\N}{\mathbb{N}}
\newcommand{\Z}{\mathbb{Z}}
\newcommand{\Q}{\mathbb{Q}}
\newcommand{\R}{\mathbb{R}}
\newcommand{\C}{\mathbb{C}}
\newcommand{\F}{\mathbb{F}}
\newcommand{\A}{\mathbb{A}}
\newcommand{\PP}{\mathbb{P}}
\newcommand{\E}{\mathbb{E}}
\newcommand{\eps}{\varepsilon}
\newcommand{\dd}{\,\mathrm{d}}
\begin{document}
\section*{UnsolvedMath ID: 20002312}
\textbf{Source problem code:} catalogue entry
\subsection{Definitions and mathematical statement}
The recovered Question 10 asks whether relative \(\kappa^+\)-computable categoricity implies relative \(\lambda^+\)-computable categoricity. Its preceding context, supplied by companion records 20002310 and 20002311, is Knight's admissible-computability section, a quasiminimal excellent class \(\mathcal K\), and cardinals \(\lambda<\kappa\). That context says \(\kappa^+\)-enumerable means \(\Sigma_1\) on \(L_{\kappa^+}\), and describes relative categoricity by isomorphisms in \(L_{\kappa^+}(M,N)\). The exact structure/diagram coding and meaning of ``in'' must be fixed; an unqualified membership claim is not a substitute for an effective graph definition.

One explicit convention for the aid's conditional work is: at an admissible height \(\rho\), effective predicates are \(\Delta_1\) over the selected admissible base, with atomic diagrams as oracle predicates when relativizing. Eligible presentations of structures of cardinality \(\rho\) and isomorphism graphs use fixed effective codes. Relative \(\rho\)-computable categoricity requires diagram-computable isomorphisms between the eligible copies. Whether all copies or only \(\rho\)-computable presentations are quantified over, admissibility with diagram predicates, and permitted parameters must be stated rather than imported silently from countable computability.

The supplied aid reports a downward-reflection criterion, not a solution of the original question. Put \(\theta=\kappa^+\), \(\eta=\lambda^+<\theta\), in its uncountable regular-cardinal setting. A high effective continuous Scott-family witness can restrict to a low one only under six stated outputs:
\begin{enumerate}
\item the needed parameters lie in the low effective base;
\item the high witness's club information reflects effectively to the low level;
\item all relevant syntax codes and parameter supports are bounded at the low level;
\item satisfaction of those formulas is absolute between the compared settings;
\item the restricted formulas cover the low tuples required by the Scott witness;
\item equal low witness types imply low automorphism-orbit equivalence by a hereditary argument that does not assume the desired low computable categoricity.
\end{enumerate}
Under a valid effective Scott-witness/categoricity correspondence in the chosen convention, the aid reports that these conditions yield relative \(\eta\)-computable categoricity. The record does not give the complete technical definitions of continuous witnesses, effective clubs or the reflection hypotheses, nor prove that the original assumptions imply them. This is a conditional strategy, not a fully instantiated general theorem.

A more concrete conditional equivalence concerns effectively basis-coordinatized structures. Suppose the class has a specified pregeometry closure, a \(\rho\)-computable reference copy with a computable basis indexed bijectively by \(\rho\), and a uniform effective procedure extending a matched pair of basis enumerations to an isomorphism. Then the aid reports that relative \(\rho\)-computable categoricity is equivalent to every eligible copy having a diagram-computable bijective basis enumeration. An isomorphism transports the reference basis; conversely the assumed extension procedure constructs an isomorphism from a selected basis. The extension algorithm and eligible-copy convention are essential extra hypotheses. The unresolved task is their downward reflection, not merely ordinary uncountable categoricity.

The appended Goncharov section heading is unrelated to Question 10 and is excluded.
\subsection{Short English statement}
Does relative computable categoricity descend from \(\kappa^+\) to \(\lambda^+\)? The supplied aid gives only a conditional Scott-witness reflection program and a basis-selection equivalence under an effective extension hypothesis. The uncountable computation and coding conventions remain part of the formulation.
\subsection{Sources}
\begin{itemize}
\item Knight, subsection 1.4, \emph{\(\kappa^+\)-computable categoricity}, Question 10, AIM \emph{Computable Stability Theory Problem Session}: \url{https://aimath.org/pastworkshops/computestabproblems.pdf}.
\item Local registry \texttt{data/unsolvedmath/problems.json}, permanent ID 20002312, code \texttt{catalogue entry}; canonical \texttt{aim-logic-notes.json}, \texttt{notes[87]}; updated 14 August 2026. Companion ID 20002310, code \texttt{AIM-LOGIC-0086}, appends the admissible-computability definitions, while ID 20002311, code \texttt{AIM-LOGIC-0087}, Question 9, supplies \(\lambda<\kappa\) and the quasiminimal excellent-class context.
\item The corrected-verified statement label records removal of the following Goncharov section heading. No detailed external theorem reference is supplied; the original literature field is empty.
\item Attempt 1: conditional result, problem status at run unknown, rigor label conditional, novelty confidence low; current status partially solved. The machine-generated aid is not peer reviewed and requires independent verification before citation.
\end{itemize}
\end{document}
